\documentclass{article}

\usepackage[utf8]{inputenc}
\usepackage[margin=1in]{geometry}
\usepackage[titletoc,title]{appendix}

\usepackage{amsmath,amsfonts,amssymb,mathtools}
\usepackage{graphicx,float}
\usepackage{booktabs}
\usepackage[ruled,vlined]{algorithm2e}
\usepackage{algorithmic}
\usepackage{biblatex}
\usepackage{cleveref}
\usepackage{indentfirst}
\setlength{\parindent}{2em}

\addbibresource{references.bib}

\title{PRML Spring 2026 Homework 3: Recurrent Neural Networks}
\author{Yifan Ding \\
        \small Code: \url{https://github.com/douhenhuidyf/PRML-Spring-2026/tree/main/hw3} \\
        \small Email: 23373493@buaa.edu.cn}
\date{April 25, 2026}

\begin{document}
\maketitle

% \begin{abstract}

% \end{abstract}


\section{Introduction}
\subsection{Problem Statement}
% The Air Quality dataset.
% This is a dataset that reports on the weather and the level of pollution each hour for five years. The data includes the date-time, the pollution called PM2.5 concentration, and the weather information including dew point, temperature, pressure, wind direction, wind speed and the cumulative number of hours of snow and rain.
% We can use this data and frame a forecasting problem where, given the weather conditions and pollution for prior hours, we forecast the pollution at the next hour.
Given the Air Quality dataset\cite{dataset}, which reports on the weather and the level of pollution each hour for five years,
we aim to forecast the PM2.5 concentration at the next hour based on the weather conditions and pollution levels from prior hours.
The dataset includes features such as date-time, dew point, temperature, pressure, wind direction, wind speed, and cumulative hours of snow and rain.
By leveraging these features, we can build a predictive model to estimate future pollution levels effectively.
\subsection{Problem Analysis}
This is a typical time series forecasting problem, where the goal is to predict future values based on historical data.
The dataset contains multiple features that can influence the target variable (PM2.5 concentration), making it a multivariate time series forecasting problem.
To address this, we can utilize Recurrent Neural Networks (RNNs) and their variants which are well-suited for capturing temporal dependencies in sequential data.

\subsection{Structure of the Paper}

\section{Methodology}
\label{sec:methodology}

\subsection{Metrics}
\label{sec:metrics}
To evaluate the performance of our forecasting model, we will use the following metrics:
\begin{enumerate}
    \item MAE (Mean Absolute Error): This metric measures the average magnitude of the errors in a set of predictions, without considering their direction. It is calculated as:
    \begin{equation}
    \text{MAE} = \frac{1}{n} \sum_{i=1}^{n} |y_i - \hat{y}_i|
    \end{equation}
    where \( y_i \) is the actual value and \( \hat{y}_i \) is the predicted value.
    \item RMSE (Root Mean Squared Error): This metric measures the square root of the average of squared differences between predicted and actual values. It is calculated as:
    \begin{equation}
    \text{RMSE} = \sqrt{\frac{1}{n} \sum_{i=1}^{n} (y_i - \hat{y}_i)^2}
    \end{equation}
    where \( y_i \) is the actual value and \( \hat{y}_i \) is the predicted value.
\end{enumerate}

\subsection{Recurrent Neural Networks}
\label{sec:rnn}
Recurrent Neural Networks (RNNs) are a class of neural networks designed to handle sequential data.
\begin{figure}[h!]
    \centering
    \includegraphics[width=0.8\textwidth]{figure/rnn.png}
    \caption{Architecture of a vanilla Recurrent Neural Network (RNN).}
    \label{fig:rnns}
\end{figure}
As shown in \cref{fig:rnns}, vanilla RNNs have a unique sequential architecture that composed of multiple identical layers of recurrent units(RNN cells).
They input data in a sequential manner and maintain a hidden state that captures information from previous time steps, endowing them suitable for tasks like time series forecasting.

However, vanilla RNNs can suffer from issues like vanishing and exploding gradients since all information is passed through the hidden state.
The longer the sequence, the more difficult it becomes for the model to learn long-term dependencies, which is a common problem in time series data.
Besides, vanilla RNNs can be computationally expensive and time-consuming as they need to forward sequentially and backpropagating through time(BPTT).
Thus vanilla RNNs can not scale up to long sequences, and we need to use more advanced architectures like Long Short Term Memory (LSTM)\cite{lstm} and Gated Recurrent Units (GRU)\cite{cho2014learningphraserepresentationsusing} to address these issues.
\subsection{Long Short Term Memory}
\label{sec:lstm}
Long Short Term Memory (LSTM) is a type of RNN architecture that is designed to overcome the vanishing gradient problem in vanilla RNNs.
LSTMs introduce a memory cell that can maintain information over long periods of time, allowing them to capture long-term dependencies in sequential data.
The architecture of an LSTM cell is shown in \cref{fig:lstm}.
\begin{figure}[H!]
    \centering
    \includegraphics[width=0.5\textwidth]{figure/lstm.png}
    \caption{Architecture of an LSTM cell.}
    \label{fig:lstm}
\end{figure}

\subsection{Gated Recurrent Units}
\label{sec:gru}
11111111111111111
\section{Experiments}
\label{sec:experiments}

\section{Conclusions}
\label{sec:conclusions}
111111111111111111111111111111

\printbibliography

% \begin{appendices}

% \section{MATLAB Functions}
% Add your important MATLAB functions here with a brief implementation explanation. This is how to make an \textbf{unordered} list:
% \begin{itemize}
%     \item \texttt{y = linspace(x1,x2,n)} returns a row vector of \texttt{n} evenly spaced points between \texttt{x1} and \texttt{x2}.
%     \item \texttt{[X,Y] = meshgrid(x,y)} returns 2-D grid coordinates based on the coordinates contained in the vectors \texttt{x} and \texttt{y}. \text{X} is a matrix where each row is a copy of \texttt{x}, and \texttt{Y} is a matrix where each column is a copy of \texttt{y}. The grid represented by the coordinates \texttt{X} and \texttt{Y} has \texttt{length(y)} rows and \texttt{length(x)} columns.
% \end{itemize}

% \end{appendices}

\end{document}
